\documentclass[10pt,a4paper]{amsart}
\usepackage[margin=19mm]{geometry}
\usepackage{amsmath,amssymb,mathtools}
\usepackage[scr=rsfs,cal=euler]{mathalfa}
\usepackage{xfrac}
\usepackage[unicode,hidelinks,bookmarks=false]{hyperref}
\newcommand{\ie}{i.e.\ }
\newcommand{\cf}{cf.\ }
\newcommand{\resp}{respectively\ }
\newcommand{\Subsection}{\subsection}
\providecommand{\nobreakdash}{\nobreak\hbox{-}}
\setlength{\emergencystretch}{2em}
\multlinegap0pt
\usepackage{bm}
\usepackage{bbm}
\usepackage{dsfont}
\usepackage{xspace}
\usepackage{cancel}
\usepackage{mathtools}
\usepackage[shortlabels]{enumitem}
\usepackage{amsthm}
\makeatletter
\@ifundefined{theorem}{\newtheorem{theorem}{Theorem}}{}
\@ifundefined{lemma}{\newtheorem{lemma}[theorem]{Lemma}}{}
\makeatother
\newcommand{\alt}{~\mid~}
\newcommand{\ccat}[1]{\mathbf{c #1}}
\newcommand{\iid}{\mathrm{id}}
\newcommand{\ovl}[1]{\overline{#1}}
\newcommand{\scat}[1]{\mathbf{\ovl{#1}}}
\newcommand{\symplus}{\gamma^\oplus}
\newcommand{\symtimes}{\gamma^\otimes}
\newcommand{\tcat}[1]{\mathbf{\ovl{#1}}_2}
\title[One rig to control them all]{One rig to control them all}
\author{Chris Heunen, Robin Kaarsgaard,, Louis Lemonnier}
\date{Source: arXiv:2510.05032v3 (CC BY 4.0)}
\begin{document}
\maketitle

\noindent\emph{Source and licence.} One rig to control them all. Version arXiv:2510.05032v3 (CC BY 4.0), distributed under the Creative Commons Attribution 4.0 International licence, as recorded in the arXiv licence metadata for this exact version and shown on the abstract page https://arxiv.org/abs/2510.05032v3. Full rights notice: licenses/arxiv-2510.05032-CC-BY-4.0.txt.\par\smallskip

\noindent\textit{Editorial scope.} Counted unit: the Lemma whose source label is \texttt{lem:b-def} in the article (source0.tex lines 1270--1272, source label \texttt{lem:b-def}), delivered together with the article's own complete original proof of it (source0.tex lines 1273--1312, one \texttt{proof} environment of 40 lines). No mathematics is written, repaired or re-derived: statement and proof are the article's own text line for line. Required context retained uncounted: cprop:complementarity (lines 717--737); eq:mon-tensor (lines 1024--1032). Every definition, display and label of those lines is retained unchanged. Omitted from this package and not redistributed: 1--716, 738--1023, 1033--1269, 1313--1865. Nothing inside a retained excerpt is omitted or altered: every retained body is delivered exactly as the article's lines are written, so no omission receipt of any kind accompanies this package. Citations: the retained excerpts carry no citation command at all, so no bibliography entry of the article is transcribed and no reference is redistributed. Reference adaptation: none was needed and none was made; every reference inside the delivered unit resolves inside it, so no unresolved reference or citation is produced. Printed slips kept literal: no printed word, symbol, hypothesis or number of the counted statement or of its proof was altered. Layout and class: the article's own class is \texttt{article} with packages including geometry, amsmath, amsthm, amssymb, inputenc, hyperref, wrapfig, anyfontsize, proof, mathrsfs, stmaryrd, dsfont, cleveref, hypcap; the wrapper is the lane's pinned \texttt{amsart} editorial wrapper at 10pt on A4, which restates the theorem-like environments the retained text needs and adds the article's own macro definitions that the retained text needs (\texttt{\textbackslash alt}, \texttt{\textbackslash ccat}, \texttt{\textbackslash iid}, \texttt{\textbackslash ovl}, \texttt{\textbackslash scat}, \texttt{\textbackslash symplus}, \texttt{\textbackslash symtimes}, \texttt{\textbackslash tcat}), with extra wrapper packages (bm, bbm, dsfont, xspace, cancel, mathtools, enumitem). The delivered numbering is the unit's own. Assets: no retained span contains a figure environment, a graphics-inclusion command, a PDF-page inclusion, a table or a TikZ picture; no image file is copied into this package, the package declares no asset dependency and no image file is redistributed with it. Licence: version arXiv:2510.05032v3 of this article is distributed under the Creative Commons Attribution 4.0 International licence, as recorded in the arXiv licence metadata for this exact version and shown on the abstract page https://arxiv.org/abs/2510.05032v3. A full rights notice is delivered at licenses/arxiv-2510.05032-CC-BY-4.0.txt. Characters: every retained excerpt is pure ASCII, so the delivered engine, XeLaTeX with its default Latin Modern text font, reports no missing character.
	\begin{enumerate}[(\alph*)]
		\item\label{cprop:composition} composition:
			%\hfill
			$C^1(g \circ f) = C^1(g) \circ C^1(f)$;
		\item\label{cprop:identity} identity:
			$C^1(\iid_n) = \iid_{n+1}$;
		\item\label{cprop:strength} strength:
			$C^1(f + \iid_m) = C^1(f) + \iid_m$;
		\item\label{cprop:colour} colour change:
			$(x + \iid_n) \circ C^0(f) \circ (x + \iid_n) = C^1(f)$;
		\item\label{cprop:complementarity} complementarity:
			$C^0(f) \circ C^1(f) = \iid_1 + f$;
		\item\label{cprop:commutativity} commutativity:
			$C^0(f_1) \circ C^1(f_2) = C^1(f_2) \circ C^0(f_1)$;
		\item\label{cprop:swap} ``swap'':
			$C^1(x) \circ \sigma_{1,1} \circ C^1(x) \circ \sigma_{1,1} \circ C^1(x)
			= \sigma_{1,1}$;
		\item\label{cprop:swapcoh} ``swap'' coherence:
			$(\sigma_{1,1} + \iid_n) \circ C^1(C^1(f)) = C^1(C^1(f))
			\circ (\sigma_{1,1} + \iid_n)$.
	\end{enumerate}

	\begin{align}
		\label{eq:mon-tensor}
		\widetilde R =~&
		\left\{ \widetilde{f + \iid} \circ \widetilde{\iid + h} = \widetilde{\iid + h} \circ
		\widetilde{f + \iid} \alt f,h \in G \right\} \\
		\label{eq:base-rel}
		&\cup
		\left\{ \widetilde f = \widetilde h \alt (f = h) \in R \right\}
	\end{align}

\begin{lemma}\label{lem:b-def}
	The functor $B \colon \tcat P \to \ccat P$ is a well-defined prop morphism.
\end{lemma}

\begin{proof}
	Note that $\scat P$ is generated by $\ovl G$ and permutations
	$\symplus$. Therefore, morphisms are of the form:
	\[
		\gamma_0 \circ (\iid \oplus \widetilde{g_1} \oplus \iid) \circ
		\gamma_1 \circ \dots \circ \gamma_{n-1} \circ (\iid \oplus
		\widetilde{g_n} \oplus \iid) \circ \gamma_n
	\]
	where $\gamma_i$ are short for compositions of symmetries for the direct
	sum $\oplus$. Since the generator terms are surrounded by permutations, we
	can write them as follows, without loss of generality:
	\[
		\gamma_0 \circ (\iid \oplus \widetilde{g_1}) \circ \gamma_1 \circ \dots
		\circ \gamma_{n-1} \circ (\iid \oplus \widetilde{g_n}) \circ \gamma_n
	\]
	If this is a morphism $2^m \to 2^m$ in $\tcat P$, then the terms containing
	generators $\widetilde g \colon 2^k \to 2^k$ are of the form:
	\[
		\iid_{2^m - 2^k} \oplus \widetilde g
		=
		\iid_{2^{m-1}} \oplus (\iid_{2^{m-2}} \oplus (\dots \oplus (\iid_{2^k} \oplus \widetilde g) \dots ))
	\]
	Therefore the definition of $B$ above defines an image of all morphism in
	$\tcat P$. We now prove that $B$ is well-defined. All coherence equations
	hold because $\beta$ is well-defined. The only equation specific to $\scat
	P_2$ is (\ref{eq:mon-tensor}), which is preserved since $\otimes$
	is a monoidal structure, provided that $B$ is indeed a prop morphism. We
	thus need to prove that $B(f \otimes g) =
	B(f) + B(g)$ for all $f$ and $g$. We know, through $\beta$, that $B(\symtimes_{2^m,2^n}) =
	\sigma_{m,n}$, so that it suffices to show that $B(\iid_2 \otimes
	f) = \iid_1 + B(f)$.
	\begin{align*}
		B(\iid_2 \otimes f)
		&= B(f \oplus f) \\
		&= B(f \oplus \iid_n) \circ B(\iid_n \oplus f) \\
		&= C^0(B(f)) \circ C^1(B(f)) \\
		&\stackrel{\ref{cprop:complementarity}}{=} \iid_1 + B(f)
	\end{align*}
	Thus $B$ is a prop morphism.
\end{proof}

\end{document}
